%File: CARMA.tex
% Ported from arr/ARRv3.tex (ACL format) into the AAAI/ICWSM template.
% Content copied as-is from the ARR submission for now; see icwsm/action-plan.md
% for the list of open issues (reviewer-raised and AAAI-formatting) still to resolve.
\documentclass[letterpaper]{article} % DO NOT CHANGE THIS
\usepackage[submission]{aaai2026}  % DO NOT CHANGE THIS
\usepackage{times}  % DO NOT CHANGE THIS
\usepackage{helvet}  % DO NOT CHANGE THIS
\usepackage{courier}  % DO NOT CHANGE THIS
\usepackage[hyphens]{url}  % DO NOT CHANGE THIS
\usepackage{graphicx} % DO NOT CHANGE THIS
\urlstyle{rm} % DO NOT CHANGE THIS
\def\UrlFont{\rm}  % DO NOT CHANGE THIS
\usepackage{natbib}  % DO NOT CHANGE THIS AND DO NOT ADD ANY OPTIONS TO IT
\usepackage{caption} % DO NOT CHANGE THIS AND DO NOT ADD ANY OPTIONS TO IT
\frenchspacing  % DO NOT CHANGE THIS
\setlength{\pdfpagewidth}{8.5in} % DO NOT CHANGE THIS
\setlength{\pdfpageheight}{11in} % DO NOT CHANGE THIS
\pdfinfo{
/TemplateVersion (2026.1)
}

% --- Additional packages needed by the ported content (not part of the
% mandatory AAAI block above). tcolorbox/xcolor/listings from the ARR
% preamble were unused in the body and have been dropped.
\usepackage{booktabs}
\usepackage{multirow}
\usepackage{microtype}

% Paper-checklist macros (from arr/paper-checklist.tex, the official AAAI ICWSM
% checklist template appended in Appendix H below).
\usepackage{xcolor}
\newcommand{\answerYes}[1]{\textcolor{blue}{#1}}
\newcommand{\answerNo}[1]{\textcolor{teal}{#1}}
\newcommand{\answerNA}[1]{\textcolor{gray}{#1}}
\newcommand{\answerTODO}[1]{\textcolor{red}{#1}}

% Arabic inline text support, carried over from ARRv3.tex (\foreignlanguage{arabic}{...}
% is used extensively in the keyword/phrase tables and worked examples).
% NOTE: AAAI camera-ready rules say non-Roman scripts must be restricted to
% bit-mapped figures -- this needs to be checked against the actual ICWSM
% author kit before camera-ready (see icwsm/action-plan.md, section 2).
% This babel/XeLaTeX-Arabic setup requires compiling with LuaLaTeX or XeLaTeX,
% same as the source ARR file.
\usepackage[english,bidi=default]{babel}
\babelfont{rm}{TeXGyreTermesX}
\babelprovide[import]{arabic}
\babelfont[*arabic]{rm}{Noto Sans Arabic}

% Section cross-references in the ported text use \S\ref{...} extensively,
% which requires numbered sections (the AAAI default of secnumdepth=0 would
% make every \ref{} print nothing). Numbering is optional for ICWSM; revisit
% before camera-ready.
\setcounter{secnumdepth}{2}

\title{CARMA: Large-Scale Arabic Reddit Corpus for Mental Health Analysis}

% AAAI anonymous-submission instructions: write "Anonymous Submission" as the
% sole author and leave affiliations empty.
\author{
    Anonymous Submission
}
\affiliations{
    ~
}

\begin{document}

\maketitle

\begin{abstract}
Mental health has emerged as a critical global issue, yet Arabic NLP resources for mental health research remain scarce. We introduce \textbf{CARMA}, a large-scale corpus from Reddit (4,365 diagnosed users and 2,875 control, \textbf{\textasciitilde625K} posts), covering users who have self-reported diagnoses of mental health conditions. The corpus targets 12 conditions and a control group. We collect data using a precision-oriented, pattern-based extraction pipeline (estimated at \textasciitilde78\% precision by manual annotation) validated
through manual human annotation that combines self-disclosure phrases with condition-specific keywords designed to support four
Arabic dialect groups (Gulf, Egyptian, Levantine, and Modern Standard Arabic) to identify posts containing genuine diagnostic self-reports, though the resulting corpus skews toward Egyptian and Gulf Arabic rather than representing these groups evenly. We employ an author-level approach that retrieves complete posting histories including relevant meta data (subreddit, time of creation, interaction stats) to enable longitudinal analysis. We benchmark the performance of current detection models and leverage Large Language Models to establish baselines for Arabic mental health classification.
\end{abstract}

\section{Introduction}
\label{sec:introduction}
Mental health has emerged as a critical global challenge, with over one billion
people worldwide estimated to live with a mental disorder
\cite{who2025mental}. In recent years, social media platforms have become a valuable source of
large-scale, naturalistic text for studying mental health-related language.
Computational approaches have demonstrated that linguistic and behavioral signals
in user-generated content can serve as proxies for detecting mental health
conditions \cite{de2013predicting, chancellor2020methods}. The CLPsych workshop
series further institutionalized this line of research by establishing datasets, shared tasks and standardized evaluation frameworks \cite{coppersmith2015clpsych}.
A central methodological development is the \textit{self-reported diagnosis}
paradigm, which identifies users who publicly disclose having received a clinical
diagnosis. First introduced by \citet{coppersmith2014quantifying} on Twitter and
later extended to Reddit by \citet{cohan-etal-2018-smhd}, this approach has enabled the
construction of large-scale datasets without requiring manual clinical annotation.

Despite this progress, the overwhelming majority of existing resources target English-language data \cite{harrigian2021state}. Arabic remains substantially underrepresented in computational mental health research, despite being spoken by over 400 million people across diverse linguistic and cultural contexts. Arabic poses particular challenges for natural language processing, including rich morphology and extensive dialectal variation across regional varieties such as Gulf, Egyptian, and Levantine Arabic \cite{darwish2021panorama,habash2010introduction}. Existing Arabic mental health datasets are limited in scale, typically target only one or two conditions (most commonly depression or suicidal ideation), and rarely address dialectal variation systematically (see Table~\ref{tab:arabic_datasets}).

In this paper, we address these gaps by introducing \textbf{CARMA}, a large-scale and carefully-curated Arabic mental
health corpus collected from Reddit. The corpus follows the self-reported
diagnosis paradigm: we develop a dialect-aware, pattern-based extraction pipeline. We adopt an author-level collection strategy, retrieving the complete posting histories of
identified users to support longitudinal analysis of linguistic patterns beyond the initial disclosure posts and including further metadata about the post. The corpus targets 12 conditions and a control group (Table~\ref{tab:arabic_datasets}).\footnote{Throughout this paper, \textit{diagnosed} and \textit{control} are shorthand for \textit{self-reported} group assignment, not clinical verification: ``diagnosed'' users are those who publicly self-disclosed a diagnosis, and ``control'' users are those screened for the absence of such disclosure, not clinically confirmed to be healthy. See \S\ref{subsec:self_disclosure_grounding} for the full discussion and its limitations.}

Our contributions are as follows:
\begin{enumerate}
    \item We introduce \textbf{CARMA}, a large-scale Arabic mental health corpus
    collected from Reddit, targeting 12 conditions and a control group.

    \item We develop the first precision-oriented, dialect aware, automatic pattern-based extraction pipeline to identify self disclosure in Arabic.

    \item We validate the automatic extraction pipeline through manual annotation.

    \item We provide a comprehensive analysis of corpus characteristics and
    baseline model performance, establishing a benchmark for Arabic mental health detection.

\end{enumerate}



\section{Related Work}
\label{sec:related_work}

\subsection{Mental Health Detection from Social Media}
\label{subsec:related_social_media}
Research on detecting mental health conditions from social media has grown substantially over the past decade, from early work demonstrating that linguistic and behavioral signals can identify users experiencing depression \citep{de2013predicting} to the CLPsych workshop series \citep{coppersmith2015clpsych}, which introduced standardized shared tasks and evaluation frameworks that helped establish methodological foundations for the field.

A key development is the \textit{self-reported diagnosis} paradigm, which identifies users who publicly disclose a clinical diagnosis through pattern matching. \citet{coppersmith2014quantifying} first introduced this approach on Twitter via regular-expression matching; \citet{cohan-etal-2018-smhd} extended it to Reddit with SMHD, a large-scale corpus covering nine conditions, and \citet{mankarious2025mindset} followed a similar methodology at larger scale on the same platform with MINDSET. The paradigm has since been applied to other languages, including German (SMHD-GER, \citealp{zanwar2023smhd}). \citet{harrigian2021state} survey existing social media mental-health datasets, noting that the majority target English, leaving other languages substantially underrepresented.

\subsection{Arabic NLP for Mental Health}
\label{subsec:related_arabic}

Compared to English, computational research on mental health in Arabic remains
relatively limited. Arabic presents several challenges for natural language
processing, including rich morphology, substantial dialectal variation across
regional varieties such as Gulf, Egyptian, and Levantine Arabic, and limited
availability of large annotated datasets \cite{darwish2021panoramic, alayba2025arabic}.
Table~\ref{tab:arabic_datasets} summarizes existing Arabic datasets related to
mental health and situates our corpus within this landscape.

Several studies have constructed X (Twitter)-based datasets targeting individual conditions: OCD symptoms \citep{al2024classification}, suicidal ideation \citep{abdulsalam2024detecting, baghdadi2022optimized}, and, most prominently, depression \citep{hassib2022aradepsu, almouzini2019detecting, elramly2021cairodep, maghraby2022mood}, including PHQ-9-organized \citep{maghraby2022mood} and comorbid depression/suicidal-ideation \citep{hassib2022aradepsu} variants. Beyond social media, \citet{alhuzali2024mentalqa} introduce MentalQA, a question-answering dataset from the Altibbi medical platform rather than user-generated discourse. Table~\ref{tab:arabic_datasets} details dataset sizes, target conditions, and annotation methods.

Despite these efforts, existing Arabic mental health datasets exhibit three
notable limitations. First, most datasets focus on only one or two conditions,
typically depression or suicidal ideation, leaving the broader spectrum of mental
health conditions largely unaddressed. Second, annotations are typically performed
at the post level, without modeling user-level context or longitudinal behavior.
Third, dialectal variation across Arabic varieties is rarely addressed
systematically during dataset construction. Our corpus is designed to address
all three of these limitations.

\begin{table*}
\centering
\small
\begin{tabular}{cl l rll}
\toprule
& \textbf{Dataset} & \textbf{Condition} & \textbf{Size} & \textbf{Domain} & \textbf{Anno.} \\
\midrule
& MentalQA~\cite{alhuzali2024mentalqa} & N/A & 1,000 & Altibbi & Human \\
& OCDTweets~\cite{al2024classification} & OCD symptoms & 5K & X & Human \\
& Arabic Suicidal Tweets~\cite{abdulsalam2024detecting} & Suicidal ideation & 5.7K & X & Human \\
& Suicide Tweets~\cite{baghdadi2022optimized} & Suicidal ideation & 2.0K & X & Human \\
& AraDepSu~\cite{hassib2022aradepsu} & Depression, Suicidal ideation & 20K & X & Human \\
& Detecting Depression~\cite{almouzini2019detecting} & Depression & 2.7K & X & Human \\
& CairoDep~\cite{elramly2021cairodep} & Depression & 7K & X & Human \\
& Mood Change~\cite{maghraby2022mood} & Depression & 49K & X & Auto. \\
\midrule
& \textbf{Ours} & \textbf{12 conditions} & \textbf{625K} & \textbf{Reddit} & \textbf{Auto.} \\

\bottomrule
\end{tabular}
\caption{Size and coverage comparison of Arabic Mental Health Corpora.}

\label{tab:arabic_datasets}
\end{table*}



\section{Dataset Construction}
\label{subsec:data_source}

\subsection{Diagnosed Dataset}

We collect data from Reddit using pre-downloaded JSONL archives from Arctic Shift\footnote{\url{https://github.com/ArthurHeitmann/arctic_shift}} (Pushshift format), containing posts and comments from Arabic-speaking subreddits\footnote{We obtained an excempt IRB approval for this study from our institution, exempt under DHHS regulatory category 4. }. Reddit offers longer, more detailed self-reports that are well-suited to the self-reported diagnosis paradigm. To the best of our knowledge, there are no dedicated Arabic-language mental health subreddits comparable to those in English (e.g., r/depression, r/SuicideWatch). Consequently, our dataset is derived entirely from general regional subreddits (e.g., r/saudiarabia, r/Egypt, r/jordan), from which we extract mental-health-relevant content; this subreddit pool is itself unevenly represented in the resulting corpus rather than uniform across regions (Appendix~\ref{app:subreddit_time}).
Our extraction pipeline follows the self-reported diagnosis paradigm established by \citet{coppersmith2014quantifying} and adopted in subsequent corpus-construction work \citep{coppersmith2015clpsych, cohan-etal-2018-smhd}, which identifies diagnosed users via precision-oriented pattern matching over diagnosis-disclosure phrases rather than clinical verification (See section \ref{subsec:self_disclosure_grounding} for more discussion about self reported diagnosis).
We use a curated list of condition keywords (e.g., \foreignlanguage{arabic}{اكتئاب حاد}, "Severe Depression") and diagnosis phrases (e.g., \foreignlanguage{arabic}{تم تشخيصي بـ}, "I was diagnosed with"), checking whether they co-occur within a 40-character window (measured at the Unicode character level), the same window width used by \citet{cohan-etal-2018-smhd} in constructing SMHD, where a sweep over candidate distances found 40 characters to maximize $F_{0.5}$ (precision-weighted) on their validation set. See Appendix~\ref{app:keywords} and Appendix~\ref{app:phrases} for a detailed pattern description. To enhance inclusivity and robustness, the phrase list incorporates variations across Arabic dialects, including Modern Standard Arabic, Egyptian, and Levantine, along with frequent spellings and misspellings to account for the diverse posting habits of social media users.


\subsection{Grounding Self-Disclosure-Based Group Assignment}
\label{subsec:self_disclosure_grounding}

Both our diagnosed and control groups are assigned entirely from users'
voluntary public self-disclosures rather than clinical verification,
following the self-reported diagnosis paradigm discussed in
\S\ref{subsec:related_social_media} \citep{coppersmith2014quantifying,
coppersmith2015clpsych, cohan-etal-2018-smhd}; \citet{de2013predicting}
establishes that such self-identified users exhibit measurable
linguistic and behavioral differences from the general population,
motivating self-disclosure as a scalable diagnosis proxy at corpus
scale.

We explicitly acknowledge the limitations this design carries. First,
self-disclosure is a proxy for diagnosis, not a diagnosis itself: users
may exaggerate or use diagnostic language loosely, genuinely diagnosed
users who never post a matching disclosure are missed (false
negatives), and sarcastic, fictional, or third-person mentions can pass
pattern matching, we estimate the resulting irregularity rate through
manual annotation (\S\ref{subsec:human_annotation}). Second, control
users are only \textit{negatively} defined by the absence of detected
mental-health language across their retrieved history, so the control
group is better understood as a screened general population than a
clinically verified healthy cohort. These limitations are inherent to
the self-reported diagnosis paradigm and are shared by the prior work
it is based on \citep{coppersmith2014quantifying, cohan-etal-2018-smhd}.

\subsection{Preprocessing}

We apply two filtering steps to the collected post histories. First, we remove all posts containing mental health-related content (self-disclosures, condition keywords, and medication mentions) \footnote{We used a list of 305 medications and 140 condition specific terms. Any post that contained an item of this list was excluded from the dataset. 14,621 posts were excluded.} to prevent models from relying on explicit diagnostic markers when learning to classify users. This design encourages models to learn more subtle linguistic and behavioral signals rather than trivial keyword-based cues, aligning with prior work in mental health NLP \cite{cohan-etal-2018-smhd}. Second, we filter out non-Arabic content using the FastText library to ensure linguistic consistency across the corpus. Any post containing 20\% or more non Arabic content was excluded. In addition, we remove posts with fewer than 10 words, including posts marked as [removed] by the Reddit API.\footnote{Arctic Shift API allows users to submit requests to remove their posts}

\subsection{Control Dataset}
We further construct a control dataset of 2,875 users in a similar manner by scraping entire public user post histories from general Arabic-speaking subreddits (e.g., r/SaudiArabia) for candidate control users. Control users are subject to strict exclusion criteria: we immediately discard a candidate if they have ever posted about a mental health-related issue on the platform. For this filtration, we use the same pattern developed to identify diagnosed users, scanning for matching phrases, medication keywords, or other mental health-related terminology. See Appendix~\ref{app:phrases} for more detail. Such stringent exclusion criteria are inspired by previous literature \cite{cohan-etal-2018-smhd} to minimize the risk of including affected users in the control group.

\section{Annotation and Validation}
\label{sec:annotation}

\subsection{Human Annotation}
\label{subsec:human_annotation}
Two native Arabic-speaking annotators
independently labeled a stratified sample of 240 posts from CARMA (20
posts per retained condition) under a binary scheme:
\texttt{TP} (a genuine first-person self-disclosure of a mental health
diagnosis) or \texttt{FP} (a clear mismatch:
sarcasm, third-person reference, figurative use, or negation). Appendix~\ref{app:guidelines}
reproduces the full written guidelines, including worked examples for
each label, given to both annotators. The two annotators achieve substantial agreement ($\kappa = 0.701$,
89.8\% raw agreement on the $n{=}235$ posts both annotators labeled \footnote{Five posts were excluded due to a missing annotation by at least one of the annotators}),
which falls in the ``substantial'' band ($0.61$--$0.80$) on the
interpretation scale of \citet{landis1977measurement} used throughout this
work (Appendix~\ref{app:agreement_details}). Per-label F1 (\%) between the two human annotators is 93.5 for TP and 76.5 for FP.


Table~\ref{tab:iaa_confusion} reports the full confusion matrix between
the two annotators on the 235 jointly-labeled posts, which we use below
to derive a corpus-wide estimate of the pipeline's residual
false-positive rate.


\paragraph{Estimated false-positive rate.} Ann-1 labeled 57/235 (24.3\%)
of the sample \texttt{FP}, while Ann-2 labeled 45/235 (19.1\%)
\texttt{FP}. We treat each annotator's FP proportion as an independent noisy
estimate of the pipeline's true false-positive rate and take their mean
as our point estimate: $(21.7\%$, with the
individual annotator rates ($19.1$--$24.3\%$) indicating the spread of
this estimate. We therefore estimate that roughly one in five posts
surfaced by the pattern-based extraction pipeline is a false positive
(sarcasm, third-person reference, figurative use, or negation), which
we report as a caveat on corpus precision (§\ref{subsec:human_annotation};
Limitations). A residual false-positive rate of this magnitude is
consistent with prior manually-validated self-report extraction
pipelines for social-media mental health corpora, which likewise rely
on human spot-checks to characterise imperfect precision rather than
achieving it by construction \citep{cohan-etal-2018-smhd,
coppersmith-etal-2014-quantifying}. We describe the pipeline as
\textit{precision-oriented} rather than \textit{high-precision} in an
unqualified sense: an estimated 78.3\% precision is a substantial
improvement over unconstrained condition-keyword matching (which would
also match third-person mentions, negations, and sarcasm with no
co-occurrence or phrase constraint at all), but it is not the near-perfect
precision that phrase ``high precision'' might otherwise imply, and we
flag the residual 21.7\% false-positive rate as a corpus-precision caveat
throughout (Limitations). We restrict this headline estimate
to the two human annotators; \S\ref{subsec:llm_judges} summarizes a
parallel experiment adding two LLM judges as additional raters and why
we do not fold them into this estimate, with the full analysis in
Appendix~\ref{app:llm_fp_estimate}.

\subsection{LLM Judges as a Robustness Check}
\label{subsec:llm_judges}

As a robustness check, we additionally had two LLM judges, Qwen2-7B-Instruct \citep{yang-etal-2024-qwen2} and Jais-2-8B-Chat \citep{anwar-etal-2025-jais2}, independently label the same 240-post sample under the identical \texttt{TP}/\texttt{FP} scheme. Agreement with the human annotators is far weaker than agreement between the two humans ($\kappa{=}0.138$--$0.231$ vs.\ $\kappa{=}0.701$; Table~\ref{tab:iaa_full}, Appendix~\ref{app:llm_fp_estimate}), and both LLMs report a markedly lower FP rate ($\approx$11--14\% vs.\ $\approx$18--23\% for the humans), recovering only 21--27\% of the posts humans labeled FP. Rather than disagreeing at random, the LLM judges default to \texttt{TP} on the same sarcasm-, negation-, and figurative-language cases that limit this pipeline's precision, consistent with known LLM limitations on Arabic dialectal sarcasm \citep{abu-farha-magdy-2020-arabic}, negation \citep{garcia-ferrero-etal-2023-not}, and leniency bias in LLM-as-judge settings \citep{jain2025beyond}. We therefore exclude the LLM judges from the headline false-positive-rate estimate of \S\ref{subsec:human_annotation}; full pairwise statistics are in Appendix~\ref{app:llm_fp_estimate}.

\section{Corpus Statistics}
\label{sec:corpus}

\subsection{Overview}
\label{subsec:corpus_overview}
The resulting corpus comprises 4,365 unique diagnosed users and their associated 625K non mental-health posts. Table~\ref{tab:corpus_summary} presents the per-condition distribution of users, posts, and post-length statistics for the 12 retained conditions.

\begin{table}[h]
\centering
\small
\begin{tabular}{lrrrr}
\toprule
\textbf{Condition} & \textbf{Users} & \textbf{Posts} & \textbf{Mean len.} & \textbf{SD} \\
\midrule
Depression          & 1,831 & 251,809 & 42.2 & 63.8 \\
Anxiety             &   644 &  90,082 & 40.4 & 59.0 \\
ADHD                &   643 & 101,107 & 41.7 & 59.2 \\
Sleep Disorder      &   357 &  60,666 & 36.5 & 47.1 \\
OCD                 &   311 &  39,413 & 45.3 & 81.1 \\
Panic Disorder      &   147 &  20,849 & 40.1 & 52.3 \\
BPD                 &    98 &  10,735 & 41.9 & 55.5 \\
PTSD                &    92 &  15,627 & 42.9 & 52.3 \\
Suicidal Ideation   &    80 &  12,266 & 46.1 & 61.3 \\
Bipolar Disorder    &    76 &   8,894 & 38.7 & 48.9 \\
Schizophrenia       &    45 &   7,703 & 61.5 & 116.4 \\
Eating Disorder     &    41 &   5,386 & 35.4 & 44.1 \\
\midrule
\textbf{Total} & \textbf{4,365} & \textbf{624,537} & \textbf{41.7} & \textbf{62.1} \\
\bottomrule
\end{tabular}
\caption{Distribution of users and posts across the 12 retained mental
health conditions in CARMA.}
\label{tab:corpus_summary}
\end{table}


Once a user's self-disclosure post is confirmed by the pattern-based extraction pipeline, we retrieve their complete posting history to enable longitudinal analysis beyond the initial self-disclosure post.

\paragraph{Comorbidity.} The pattern-based extraction pipeline matches self-disclosure posts independently per condition,
so a user can, prior to filtering, be matched by more than one condition's phrase-keyword pattern.
394 users (8.3\%) have self-disclosure matches for two or more conditions.
The most frequent comorbid pairs are depression+anxiety (86 users), depression+ADHD (80), depression+OCD (40), ADHD+anxiety (33), and depression+sleep disorder (32), broadly consistent with known comorbidity patterns in the clinical literature. None of these 394 comorbid users are present in the final retained corpus of Table~\ref{tab:corpus_summary}: rather than assigning a comorbid user to a single primary condition, the classification pipeline (§\ref{sec:experiments}) excludes comorbid users entirely, so that each condition's positive class is free of other-condition self-disclosures. This design isolates condition-specific linguistic signals and avoids label ambiguity, providing a cleaner benchmark for single-condition classification; modeling comorbidity is an important direction for future work.

\subsection{Language and Dialect}
\label{subsec:dialect}
Dialect identification was performed using CAMeLBERT-Mix DID NADI Model\footnote{\url{https://huggingface.co/CAMeL-Lab/bert-base-arabic-camelbert-mix-did-nadi}} \cite{inoue-etal-2021-interplay}, a BERT-based transformer fine-tuned on the 21-label NADI country-level dataset for Arabic dialect classification; \citet{inoue-etal-2021-interplay} report this model among the strongest of the DID variants they evaluate on the NADI country-level benchmark. We did not separately validate the model's accuracy on CARMA's own Reddit text, so some domain shift relative to its original (Twitter-sourced) NADI training/evaluation data is possible; we treat the resulting dialect labels as a description of the corpus rather than a validated ground truth, consistent with how the extraction pipeline's own precision is treated (\S\ref{subsec:human_annotation}). As shown in Figure~\ref{fig:dialect_dist},
Egyptian Arabic is the dominant dialect in CARMA (56.1\%), with Egypt as the primary contributing country, consistent with its status as one of the most populous Arabic-speaking countries with approximately 116 million inhabitants as of 2024\footnote{Population estimates sourced from official 2024 UN demographic data.} \citep{un2024wpp_data}. Gulf Arabic constitutes the second-largest group, primarily driven by Saudi Arabia (approximately 35 million inhabitants). Levantine and Maghrebi Arabic together account for under 7\% of the corpus. This distribution reflects the strong presence of Egyptian Arabic speakers on Reddit and their engagement in public regional subreddits, and is corroborated directly by the subreddit-level provenance of the corpus's self-disclosure posts (Appendix~\ref{app:subreddit_time}, Figure~\ref{fig:subreddit_dist}): although the extraction pipeline's phrase and keyword patterns are designed to support Gulf, Egyptian, Levantine, and Modern Standard Arabic phrasing (\S\ref{subsec:data_source}), the resulting corpus is not evenly distributed across these dialects, and should not be read as a general or balanced sample of Arabic Reddit.


\begin{figure}[h]
    \centering
    \includegraphics[width=\linewidth]{latex/figures/fig_dialect_overall.png}
    \caption{Arabic Dialect Distribution in CARMA}
    \label{fig:dialect_dist}
\end{figure}

\subsection{Lexical-Semantic Differences}
\label{subsec:lexsem}

To complement the dialect distribution above, we characterize lexical differences between diagnosed and control users using a log-odds-ratio with an informative Dirichlet prior \citep{monroe-etal-2008-fightin} over unigrams (pooled diagnosed vs.\ pooled control corpus, combined word counts as background prior), which surfaces words distinctively over-used by one group while correcting for raw frequency (a z-score, not a raw rate difference).
To keep the ranked vocabulary content-bearing and free of single-user artifacts,
we (i) exclude a list of 2,001 stopwords, (ii) require a token to occur at least 10 times in the combined corpus, and
(iii) require that, in whichever group a token has nonzero count, at least $k{=}5$ distinct users contribute to that count.
Without this last adjustment, a handful of repetitive or high-volume posters can register as a spurious group-level
``keyword'', e.g., one control user alone accounted for all 190 occurrences of a proper noun that otherwise ranked
among the top control-associated terms.

Table~\ref{tab:lexsem_terms} lists the most distinctive terms per group after these adjustments. Diagnosed-associated terms still include Gulf-Arabic discourse markers and free pronouns/particles (e.g., \foreignlanguage{arabic}{احس} ``I feel,'' \foreignlanguage{arabic}{مدري} ``I don't know,'' \foreignlanguage{arabic}{ربي} ``my Lord''), each used by hundreds to thousands of distinct diagnosed users, alongside content words such as \foreignlanguage{arabic}{الصدق} ``the truth/honesty'' and \foreignlanguage{arabic}{اشياء} ``things.'' Control-associated terms skew toward proper nouns and topical terms (e.g., \foreignlanguage{arabic}{الجيش} ``the army,'' \foreignlanguage{arabic}{سوريا} ``Syria,'' \foreignlanguage{arabic}{منسف} ``Mansaf,'' a regional dish), consistent with the control corpus drawing more heavily from news/political and regional-lifestyle subreddits. We report this ranking as a description of the corpus rather than as validated clinical signal, since it appears to reflect dialect and topic composition at least as much as condition-specific language. This is not only a caveat on the log-odds ranking itself: the same dialect and topic composition differences plausibly confound the diagnosed-vs-control classification results reported in §\ref{sec:results}, which we discuss further in Limitations.


\begin{table}[h]
\centering
\small
\begin{tabular}{llrrr}
\toprule
\textbf{Term} & \textbf{Gloss} & \textbf{z} & $n_{D}$ & $n_{C}$ \\
\midrule
\multicolumn{5}{l}{\textit{Diagnosed-associated}} \\
\foreignlanguage{arabic}{احس} & I feel & 15.11 & 2076 & 542 \\
\foreignlanguage{arabic}{مدري} & I don't know (Gulf) & 10.58 & 1062 & 220 \\
\foreignlanguage{arabic}{ربي} & my Lord & 10.43 & 1213 & 282 \\
\foreignlanguage{arabic}{يارب} & oh God & 9.99 & 1446 & 431 \\
\foreignlanguage{arabic}{دايم} & always & 9.61 & 872 & 191 \\
\foreignlanguage{arabic}{تبي} & you want (Gulf) & 8.04 & 846 & 173 \\
\foreignlanguage{arabic}{ادري} & I know (Gulf) & 7.87 & 1281 & 346 \\
\foreignlanguage{arabic}{الصدق} & the truth / honesty & 7.81 & 649 & 141 \\
\midrule
\multicolumn{5}{l}{\textit{Control-associated}} \\
\foreignlanguage{arabic}{الجيش} & the army & $-$10.56 & 788 & 459 \\
\foreignlanguage{arabic}{المجاز} & metaphor & $-$10.54 & 15 & 9 \\
\foreignlanguage{arabic}{متفرقه} & miscellaneous & $-$9.44 & 63 & 24 \\
\foreignlanguage{arabic}{بجامعه} & at a university & $-$9.43 & 166 & 55 \\
\foreignlanguage{arabic}{سوريا} & Syria & $-$8.81 & 526 & 306 \\
\foreignlanguage{arabic}{الهدم} & demolition & $-$8.46 & 19 & 7 \\
\foreignlanguage{arabic}{منسف} & Mansaf (regional dish) & $-$8.45 & 35 & 20 \\
\foreignlanguage{arabic}{تريد} & want (MSA) & $-$8.40 & 212 & 119 \\
\bottomrule
\end{tabular}
\caption{Most distinctive unigrams per group by log-odds-ratio with Dirichlet prior.}
\label{tab:lexsem_terms}
\end{table}

\section{Experiments}
\label{sec:experiments}

We conducted two binary classification experiments to evaluate the dataset's utility: classical models trained on TF-IDF bag-of-words features, and pre-trained language models (PLMs) fine-tuned on user text. Both pipelines are evaluated at the \textit{post level} (each post an independent instance, label inherited from its author) and \textit{user level} (a user's posts aggregated into a single TF-IDF document, or a single truncated token sequence for PLMs), letting us compare how classical and PLM-based classifiers respond to the same shift in input granularity. User-level aggregation trades a much larger number of noisier training instances for fewer, richer per-user representations, but also shrinks the effective test-set size for lower-prevalence conditions (Appendix Tables~\ref{tab:results_classical_full} and~\ref{tab:results_plm_full}). Both experiments use a balanced 1:1 ratio of diagnosed to control samples.

\subsection{Experimental Setup}
\label{subsec:setup}

We frame the task as binary classification: for each condition $C$, classify whether a user is diagnosed with $C$ (positive) or is a control user (negative), randomly sampled from Arabic-speaking Reddit users with no self-reported mental health conditions. We split the data into 90\% training and 10\% test sets, ensuring no user appears across multiple splits; for post-level experiments, posts are grouped by user prior to splitting to prevent leakage.


\subsection{Models}
\label{subsec:models}

\textbf{Classical Classifiers.} We train four traditional machine learning models on TF-IDF bag-of-words features: Naive Bayes (NB), Logistic Regression (LR), Support Vector Machines (SVM), and XGBoost (XGB). These serve as baselines to assess whether simple lexical features can capture condition-specific language patterns.

\textbf{Pre-trained Language Models (PLMs).} We fine-tune two Arabic PLMs:
\begin{itemize}
    \item \textbf{AraBERT-Twitter} (\texttt{aubmindlab/bert-base-arabertv02-twitter}): A BERT model pre-trained on Arabic tweets, well-suited for social media text.
    \item \textbf{CAMeLBERT} (\texttt{CAMeL-Lab/bert-base-arabic-camelbert-mix}): A BERT model pre-trained on a diverse mix of Arabic text including Modern Standard Arabic and dialectal content.
\end{itemize}
For fine-tuning, we use a learning rate of 2e-5, batch size of 16, and train for 3 epochs with early stopping based on validation F1.

\subsection{Evaluation}
\label{subsec:eval}

We evaluate all models using four standard metrics: F1-score, Accuracy, Precision, and Recall. F1-score serves as our primary metric as it balances precision and recall, which is important given the class imbalance across conditions. All metrics are computed on the held-out test set.


\section{Results and Discussion}
\label{sec:results}
In this work, we introduce a large-scale Arabic mental health corpus comprising over 625K posts spanning 12 diagnosed mental health conditions. The primary objective of this resource is to facilitate future research in Arabic mental health analysis by providing a publicly available benchmark, rather than to establish new state-of-the-art performance on mental health classification.

Table~\ref{tab:results_summary} summarizes the best F1 score achieved for each condition by the two modeling pipelines at both the post and user levels (reported as \textit{post level $|$ user level}). Complete per-model results are provided in Appendix Tables~\ref{tab:results_classical_full} and~\ref{tab:results_plm_full}.

Four considerations are important when interpreting these results. First, the reported scores should be viewed primarily as strong baseline benchmarks for future work rather than as an exhaustive evaluation of the full range of possible modeling approaches. Second, performance on conditions with relatively few diagnosed users (e.g., eating disorders, schizophrenia, bipolar disorder, and suicidal ideation) should be interpreted with caution, as the corresponding test sets are small and therefore produce less stable evaluation estimates. Third, and most importantly, dialectal and topical differences between the diagnosed and control subreddit pools are a plausible confound that neither pipeline can separate from condition-specific signal at the representation level, so the F1 scores below should be read as evidence of separability under the current corpus construction, not as validated clinical signal; see §\ref{subsec:lexsem} and Limitations for the supporting evidence and full discussion. Fourth, the notably high recall scores at the post level (Appendix Table~\ref{tab:results_classical_full}) partly reflect a post-count imbalance introduced by our user-level balancing: diagnosed and control users are sampled 1:1 by \textit{user count} (\S\ref{subsec:setup}), but diagnosed users contribute roughly 1.7--2.6$\times$ as many posts per user on average as control users across conditions, so post-level train/test sets skew toward the positive class even though user counts are balanced, biasing classifiers toward predicting positive and inflating recall (at some cost to precision) relative to what a post-count-balanced split would show.

\begin{table}[t]
\centering
\small
\begin{tabular}{lcc}
\toprule
\textbf{Condition} & \textbf{Classical} & \textbf{PLM} \\
\midrule
Depression & .81$|$.68 & .80$|$.71 \\
Anxiety & .83$|$.85 & .80$|$.72 \\
ADHD & .75$|$.72 & .75$|$.71 \\
Sleep dis. & .81$|$.89 & .78$|$.73 \\
OCD & .80$|$.82 & .71$|$.72 \\
Panic & .82$|$.91 & .85$|$.77 \\
BPD & .86$|$.72 & .85$|$.71 \\
PTSD & .67$|$.70 & .70$|$.64 \\
Suicidal & .86$|$.82 & .87$|$.78 \\
Bipolar & .88$|$.59 & .86$|$.71 \\
Schizophrenia & .96$|$.91 & .88$|$.80 \\
Eating dis. & .79$|$.73 & .75$|$.73 \\
\bottomrule
\end{tabular}
\caption{Best F1 score per condition on CARMA (Reddit). Each cell reports \textit{post level $|$ user level}. Full results in Appendix Tables~\ref{tab:results_classical_full} and~\ref{tab:results_plm_full}.}
\label{tab:results_summary}
\end{table}

\subsection{Classical Classifier Performance}

Schizophrenia achieves the highest post-level F1 (.956, NB); PTSD is hardest (.665). XGBoost attains the highest post-level mean F1 (.815), narrowly ahead of NB (.812), LR (.798), and SVM (.764). Aggregating to the user level reshuffles both scores and model ranking: SVM becomes strongest on average (.730) despite being weakest at the post level, while NB drops from strongest to weakest (.644), plausibly because NB's frequent-lexical-marker strength dilutes once posts are merged, whereas SVM's margin-based boundary is more robust to the resulting feature sparsity. User-level test sets are small for several conditions (as few as 10 users; Appendix Table~\ref{tab:results_classical_full}), so individual values should be read with this variance in mind.

\subsection{Fine-tuned PLM Performance}

AraBERT-Twitter and CAMeLBERT perform comparably at both the post level (mean F1 .792 vs.\ .788) and user level (.702 vs.\ .700), with neither architecture consistently dominating (AraBERT ahead on 7/12 conditions post-level, 5/12 user-level). Both models drop substantially from post to user level, reflecting information loss from truncating a user's full history to a single fixed-length input. Schizophrenia is easiest (.875 post, .800 user) and PTSD hardest (.677--.701 post, .588--.640 user) for both models.

\subsection{Post-level vs.\ User-level Aggregation}

Aggregation reduced mean F1 for every model (classical: NB $-$.168, LR $-$.115, SVM $-$.034, XGB $-$.127; PLMs: AraBERT $-$.090, CAMeLBERT $-$.088), reflecting harder user-level representation (vocabulary dilution, input truncation) and smaller test sets for low-prevalence conditions. SVM was the most robust model to this shift, NB the most brittle. Five classical-pipeline conditions (anxiety, sleep disorder, OCD, panic, PTSD) improved at the user level, versus only a marginal OCD improvement for PLMs, suggesting TF-IDF benefits more from aggregated histories than fine-tuned PLMs.


\section{Conclusion}
We introduced CARMA, a large-scale Arabic mental health corpus comprising 4,365 diagnosed users across 12 retained mental health conditions and 2,875 control users from Reddit. Our work addresses a critical gap in Arabic NLP resources for mental health research. We developed a dialect-aware, pattern-based extraction pipeline designed to support self-reported diagnoses across Gulf, Egyptian, Levantine, and Modern Standard Arabic, though the resulting corpus is dominated by Egyptian and Gulf Arabic rather than evenly distributed across dialects (§\ref{subsec:dialect}); it was validated through manual annotation by two native Arabic-speaking annotators achieving substantial inter-annotator agreement.

Our baseline experiments, run at both the post level and the user level, demonstrate that both classical classifiers and fine-tuned pre-trained language models can distinguish diagnosed users from controls across nearly all 12 conditions (F1 up to 0.96, schizophrenia), though this separability is confounded with dialect and topic composition (§\ref{sec:results}; Limitations) and should not be read as validated clinical signal. Both pipelines lose performance on average moving from post-level to user-level aggregation, with classical models more sensitive to this shift than the PLMs.

We release this corpus to facilitate future research in Arabic computational psychiatry, dialectal NLP, and cross-cultural mental health analysis. We hope CARMA will serve as a foundation for developing more inclusive mental health technologies for the Arabic-speaking world.




\section*{Limitations}
\paragraph{Dialectal and Topical Confounding in Classification Results.}
A key concern with the classification results of §\ref{sec:results} is that reported performance may partially reflect dialectal and topical differences between the diagnosed and control groups rather than condition-specific language. The log-odds analysis of §\ref{subsec:lexsem} indicates that many of the most discriminative features are regional discourse markers (e.g., Gulf-Arabic pronouns and particles) or topic-specific vocabulary (e.g., news/political terms, regional dishes) rather than clinical language, and this pattern is consistent with the regional and dialectal skew documented in §\ref{subsec:dialect} and Appendix~\ref{app:subreddit_time}: diagnosed and control users are drawn from overlapping but not identical pools of subreddits, so dialect and topic composition can differ systematically between the two groups by construction, independent of the target condition. Neither the TF-IDF nor the fine-tuned PLM pipeline is able to separate condition-specific signal from these confounds at the representation level, and we did not run experiments that explicitly control for dialect or topic (e.g., dialect-matched control sampling, or ablations that remove the log-odds-flagged terms). Frequency- or subreddit-matched control sampling, as used in constructing SMHD \citep{cohan-etal-2018-smhd} and MINDSET \citep{mankarious2025mindset}, was feasible in those corpora because English-language Reddit offers a large enough pool of candidate control users and matched subreddits to sample from; the Arabic-speaking subreddit ecosystem is comparatively small (\S\ref{subsec:data_source}, Appendix~\ref{app:subreddit_time}), so restricting control sampling to subreddits or posting-frequency bands that match the diagnosed group would substantially shrink an already modest control pool and is not currently feasible at this corpus's scale. We therefore cannot currently quantify how much of the reported classification performance is attributable to genuine condition-specific signal versus dialectal or topical separability, and recommend the classification results in §\ref{sec:results} be interpreted as evidence of separability between diagnosed and control text under the current corpus construction, rather than as validated clinical signal, until controlled analyses of this kind are conducted. Designing and validating such controls, e.g., dialect-matched control sampling, or topic-adjusted feature ablations, is a substantial undertaking in its own right and is out of scope for the present work, whose primary contribution is the construction, validation, and release of the CARMA corpus itself rather than a definitive classification benchmark; we leave this confound analysis to future work built on the released resource.

\paragraph{Annotation Pipeline Reliability.}
The pattern-based extraction pipeline prioritises precision through strict phrase-keyword co-occurrence matching (Appendix~\ref{app:keywords}, Appendix~\ref{app:phrases}), which may miss self-disclosures phrased outside the curated phrase and keyword lists, limiting recall. Our human validation on a 240-post sample indicates substantial but imperfect precision (Ann-1 vs.\ Ann-2 $\kappa = 0.701$, 89.8\% agreement; §\ref{subsec:human_annotation}), from which we estimate a residual false-positive rate of approximately 21.7\% (mean of the two annotators' individual FP rates, 19.1--24.3\%; §\ref{subsec:human_annotation}) in the corpus.

\paragraph{Temporal Validity of Posting History.}
Once a user's disclosure post is confirmed, we retrieve and use their \textit{entire} available posting history (\S\ref{subsec:corpus_overview}) without regard to when each post falls relative to the disclosure event: a post written years before the disclosure -- potentially before the condition's onset or before the user's own diagnosis -- is treated identically to one written immediately after it. Because the released per-post fields do not retain a timestamp for every post in a user's complete history (only the confirming self-disclosure posts retain \texttt{created\_utc}; Appendix~\ref{app:subreddit_time}), we cannot directly quantify what fraction of each user's retained posts precede versus follow their disclosure. This is a limitation inherited from the underlying self-reported-diagnosis paradigm (\S\ref{subsec:self_disclosure_grounding}) rather than one specific to CARMA, but it means linguistic patterns attributed to a condition may in part reflect a user's language at any point in their posting history, not specifically the period during or after diagnosis.

\paragraph{Input Truncation in PLM Baselines.}
The user-level PLM experiments (\S\ref{subsec:setup}) concatenate a user's posts into a single sequence truncated to the model's maximum input length, discarding whatever text falls beyond that limit and collapsing the temporal order of a user's posting history into one static representation. This is a likely contributor to the post-to-user performance drop observed for both AraBERT-Twitter and CAMeLBERT (§\ref{sec:results}) and means the user-level PLM results should be read as a baseline for a single-sequence encoding of user history, not as an upper bound on what a history-aware architecture (e.g., hierarchical or windowed encoding of a user's posts) could achieve on this corpus.

\paragraph{Regional and Dialectal Skew.}
Although the extraction pipeline's phrase and keyword patterns are designed to support Gulf, Egyptian, Levantine, and Modern Standard Arabic, the resulting corpus is not a balanced sample across these varieties: Egyptian and Gulf Arabic together dominate, while Levantine and Maghrebi Arabic combined account for under 7\% of the corpus, driven by a small number of Egyptian- and Saudi-focused subreddits contributing the majority of confirmed self-disclosure posts (§\ref{subsec:dialect}, Appendix~\ref{app:subreddit_time}). Posting volume is also heavily concentrated in recent years (2024 onward; Appendix~\ref{app:subreddit_time}, Figure~\ref{fig:posts_over_time}). Findings and benchmark results should therefore not be read as representative of Arabic Reddit, or of Arabic-speaking communities, in general.

\paragraph{Sensitive Data Handling.}
\label{par:sensitive_data}
This work involves sensitive personal health information derived from social media; see Ethical Considerations below for how we mitigate re-identification and misuse risks.


\section*{Ethical Considerations}
\label{sec:ethics}
This study involves publicly available social media data from users who have voluntarily disclosed their mental health status. We acknowledge the sensitive nature of mental health data and received IRB approval from our institution in regards to the collection and distribution of this data.

All data collection adheres to Reddit's Terms of Service and API usage guidelines at the time of collection. We use pre-downloaded public archives from Arctic Shift, which contain posts from publicly accessible Arabic-speaking subreddits, collecting only content that users have chosen to share publicly.

Despite our best efforts to anonymize data, residual re-identification risks may remain. The released corpus consists of post-level text and per-user groupings with Reddit usernames replaced by irreversibly hashed pseudonymous identifiers, with all self-disclosure posts removed from each user's released history (\S\ref{par:sensitive_data}). Consistent with best practices in social media mental health research \cite{harrigian2021state}, the dataset is distributed exclusively under a Data Use Agreement (DUA) that restricts access to qualified researchers, prohibits attempts to re-identify individuals or re-link hashed identifiers to Reddit accounts, and prohibits redistribution.

The corpus is intended solely for research purposes, specifically to advance Arabic NLP for mental health detection. We recognize that users who disclose mental health conditions on social media may not anticipate their posts being used in research. We mitigate this concern by removing all self-disclosure posts from users' released posting histories, ensuring that the data cannot be used to re-identify individuals based on their diagnostic disclosures. Researchers using this corpus should adhere to ethical guidelines for mental health research and avoid applications that could stigmatize or harm the individuals represented in the data.

\paragraph{Dual-Use Risk.} A resource of this kind carries a dual-use risk beyond the individual re-identification risks addressed above: the same detection methodology could in principle be repurposed for non-consensual, individual-level mental-health profiling or screening at scale (e.g., by employers, insurers, educational institutions, or state actors) rather than for the aggregate, corpus-level research CARMA is intended for. We consider this risk to be elevated, not merely generic, in the specific social and legal context this corpus covers: across several of the jurisdictions the corpus's subreddits are drawn from, disclosed mental illness can carry significant social stigma and, in some cases, legal or occupational consequences (e.g., around employment, licensing, or family law), consequences that a self-disclosed, non-clinically-verified label (\S\ref{subsec:self_disclosure_grounding}) makes especially inappropriate to act on for any individual. We do not release any tooling packaged for individual-level screening, gate the corpus behind a DUA that restricts use to qualified researchers and prohibits re-identification (above), and explicitly discourage any deployment of models trained on this corpus for individual diagnostic, screening, employment, insurance, or law-enforcement use; the corpus and any derived models are intended for aggregate, population- and dataset-level research only.

% Bibliography entries for the entire Anthology, followed by custom entries
\bibliography{custom}

\section{Paper Checklist}
\label{sec:checklist}

\begin{enumerate}

\item For most authors...
\begin{enumerate}
    \item  Would answering this research question advance science without violating social contracts, such as violating privacy norms, perpetuating unfair profiling, exacerbating the socio-economic divide, or implying disrespect to societies or cultures?
    \answerYes{Yes, and: this work fills a documented gap in Arabic-language mental-health NLP resources (\S\ref{sec:introduction}), and we engage directly with the social-contract risks it raises -- misleading terminology, dual-use, and stigma -- in Ethical Considerations.}
  \item Do your main claims in the abstract and introduction accurately reflect the paper's contributions and scope?
    \answerYes{Yes, see the Abstract and the contributions list in \S\ref{sec:introduction}.}
   \item Do you clarify how the proposed methodological approach is appropriate for the claims made?
    \answerYes{Yes, see \S\ref{subsec:data_source}--\ref{sec:annotation}, where the extraction pipeline is described and its precision is empirically validated via manual annotation rather than assumed.}
   \item Do you clarify what are possible artifacts in the data used, given population-specific distributions?
    \answerYes{Yes, see \S\ref{subsec:dialect} and the Limitations section (Dialectal and Topical Confounding; Regional and Dialectal Skew).}
  \item Did you describe the limitations of your work?
    \answerYes{Yes, see Limitations.}
  \item Did you discuss any potential negative societal impacts of your work?
    \answerYes{Yes, see Ethical Considerations, including the Dual-Use Risk paragraph.}
      \item Did you discuss any potential misuse of your work?
    \answerYes{Yes, see the Dual-Use Risk paragraph in Ethical Considerations.}
    \item Did you describe steps taken to prevent or mitigate potential negative outcomes of the research, such as data and model documentation, data anonymization, responsible release, access control, and the reproducibility of findings?
    \answerYes{Yes, see Ethical Considerations: IRB approval, hashed pseudonymous identifiers, removal of self-disclosure posts, and Data Use Agreement (DUA) gated release.}
  \item Have you read the ethics review guidelines and ensured that your paper conforms to them?
    \answerYes{Yes.}
\end{enumerate}

\item Additionally, if your study involves hypotheses testing...
\begin{enumerate}
  \item Did you clearly state the assumptions underlying all theoretical results?
    \answerNA{NA: this paper does not test formal social-science hypotheses; it introduces a corpus and baseline classification results.}
  \item Have you provided justifications for all theoretical results?
    \answerNA{NA.}
  \item Did you discuss competing hypotheses or theories that might challenge or complement your theoretical results?
    \answerNA{NA.}
  \item Have you considered alternative mechanisms or explanations that might account for the same outcomes observed in your study?
    \answerNA{NA.}
  \item Did you address potential biases or limitations in your theoretical framework?
    \answerNA{NA.}
  \item Have you related your theoretical results to the existing literature in social science?
    \answerNA{NA.}
  \item Did you discuss the implications of your theoretical results for policy, practice, or further research in the social science domain?
    \answerNA{NA.}
\end{enumerate}

\item Additionally, if you are including theoretical proofs...
\begin{enumerate}
  \item Did you state the full set of assumptions of all theoretical results?
    \answerNA{NA: this paper contains no theoretical proofs.}
	\item Did you include complete proofs of all theoretical results?
    \answerNA{NA.}
\end{enumerate}

\item Additionally, if you ran machine learning experiments...
\begin{enumerate}
  \item Did you include the code, data, and instructions needed to reproduce the main experimental results (either in the supplemental material or as a URL)?
    \answerYes{Yes, and: the corpus is released under a Data Use Agreement restricting access to qualified researchers (Ethical Considerations); training and evaluation instructions are given in \S\ref{subsec:setup}--\ref{subsec:models}.}
  \item Did you specify all the training details (e.g., data splits, hyperparameters, how they were chosen)?
    \answerYes{Yes, see \S\ref{subsec:setup} (90/10 split, user-level grouping to prevent leakage) and \S\ref{subsec:models} (learning rate, batch size, epochs, early stopping).}
     \item Did you report error bars (e.g., with respect to the random seed after running experiments multiple times)?
    \answerNo{No, results are reported for a single run rather than averaged across seeds; we do note where small test-set sizes make individual conditions' estimates higher-variance (\S\ref{sec:results}, Appendix Tables~\ref{tab:results_classical_full}--\ref{tab:results_plm_full}).}
	\item Did you include the total amount of compute and the type of resources used (e.g., type of GPUs, internal cluster, or cloud provider)?
    \answerNo{No, compute resources are not itemized in this submission.}
     \item Do you justify how the proposed evaluation is sufficient and appropriate to the claims made?
    \answerYes{Yes, see \S\ref{subsec:eval} for the choice of metrics, and \S\ref{sec:results} for the interpretive caveats (small test sets, dialectal/topical confounding, post-level class skew) that scope what the results do and do not support.}
     \item Do you discuss what is ``the cost`` of misclassification and fault (in)tolerance?
    \answerNo{No formal cost analysis is included; the Dual-Use Risk paragraph in Ethical Considerations addresses related harms from acting on model outputs at the individual level, but does not quantify misclassification cost.}

\end{enumerate}

\item Additionally, if you are using existing assets (e.g., code, data, models) or curating/releasing new assets, \textbf{without compromising anonymity}...
\begin{enumerate}
  \item If your work uses existing assets, did you cite the creators?
    \answerYes{Yes, e.g. Arctic Shift (\S\ref{subsec:data_source}), AraBERT-Twitter and CAMeLBERT (\S\ref{subsec:models}), Qwen2-7B-Instruct and Jais-2-8B-Chat (\S\ref{subsec:llm_judges}), and the CAMeLBERT-Mix DID NADI dialect model (\S\ref{subsec:dialect}).}
  \item Did you mention the license of the assets?
    \answerNo{No, the licenses of the pretrained models and archived data source used are not enumerated in this submission.}
  \item Did you include any new assets in the supplemental material or as a URL?
    \answerYes{Yes, and: the CARMA corpus itself is the new asset, released under a Data Use Agreement (Ethical Considerations).}
  \item Did you discuss whether and how consent was obtained from people whose data you're using/curating?
    \answerYes{Yes, and: no individual consent was obtained given the public nature of the data; this is addressed via IRB exemption (DHHS category 4) and mitigated through DUA-gated release and removal of self-disclosure posts (Ethical Considerations).}
  \item Did you discuss whether the data you are using/curating contains personally identifiable information or offensive content?
    \answerYes{Yes, see Ethical Considerations (sensitive mental-health information, re-identification risk) and \S\ref{par:sensitive_data}.}
\item If you are curating or releasing new datasets, did you discuss how you intend to make your datasets FAIR (see \citet{fair})?
\answerNo{No, the FAIR framework is not discussed explicitly, though the DUA-gated release (Ethical Considerations) addresses accessibility to qualified researchers.}
\item If you are curating or releasing new datasets, did you create a Datasheet for the Dataset (see \citet{gebru2021datasheets})?
\answerNo{No, a formal Datasheet is not included in this submission.}
\end{enumerate}

\item Additionally, if you used crowdsourcing or conducted research with human subjects, \textbf{without compromising anonymity}...
\begin{enumerate}
  \item Did you include the full text of instructions given to participants and screenshots?
    \answerYes{Yes, and: Appendix~\ref{app:guidelines} reproduces the full annotation guidelines given to annotators; no screenshots are included, as annotation was not conducted through a crowdsourcing platform interface.}
  \item Did you describe any potential participant risks, with mentions of Institutional Review Board (IRB) approvals?
    \answerYes{Yes, see Ethical Considerations (exempt IRB approval, DHHS regulatory category 4).}
  \item Did you include the estimated hourly wage paid to participants and the total amount spent on participant compensation?
    \answerNo{No, annotator compensation is not described in this submission.}
   \item Did you discuss how data is stored, shared, and deidentified?
   \answerYes{Yes, see Ethical Considerations (hashed pseudonymous identifiers, removal of self-disclosure posts, DUA-gated distribution).}
\end{enumerate}

\end{enumerate}

\appendix

\section{Subreddit and Temporal Distribution}
\label{app:subreddit_time}

The released post data does not
retain a subreddit or timestamp field for every post in a user's
complete retrieved history, so a per-post breakdown over the full
$\sim$625K-post corpus is not available. We can, however, characterize
the provenance of the posts that \textit{seeded} each user's inclusion
in CARMA: the self-disclosure posts confirmed by the pattern-based
extraction pipeline (\S\ref{subsec:data_source}) do retain subreddit
and \texttt{created\_utc} metadata. Figure~\ref{fig:subreddit_dist}
plots the subreddit distribution of these 6,457 confirmed
self-disclosure posts (across the 12 retained conditions), and
Figure~\ref{fig:posts_over_time} plots their monthly volume over time.

\begin{figure}[h]
    \centering
    \includegraphics[width=\linewidth]{latex/figures/fig_subreddit_dist.png}
    \caption{Subreddit distribution of the self-disclosure posts confirming diagnosed users in CARMA. The top 5 subreddits account for roughly 69\% of confirmed disclosures; the remaining 35 subreddits contribute the rest.}
    \label{fig:subreddit_dist}
\end{figure}

\begin{figure}[h]
    \centering
    \includegraphics[width=\linewidth]{latex/figures/fig_posts_over_time.png}
    \caption{Monthly volume of self-disclosure posts confirming diagnosed users in CARMA, 2018--2026. Volume is negligible before 2022 and rises sharply from 2024 onward.}
    \label{fig:posts_over_time}
\end{figure}

Two observations follow directly from these figures and are referenced
in \S\ref{subsec:data_source} and \S\ref{subsec:dialect}. First, subreddit
coverage is highly concentrated: the five largest subreddits
(predominantly Egyptian- and Saudi-focused: \textit{r/AlexandriaEgy},
\textit{r/saudiarabia}, \textit{r/CAIRO}, \textit{r/SaudiForSaudis},
\textit{r/ExEgypt}) account for roughly 69\% of confirmed
self-disclosure posts, while Levantine-associated subreddits
(e.g., \textit{r/jordan}) contribute under 6\% and no clearly
Maghrebi-associated subreddit appears among the top contributors. This
subreddit concentration is consistent with, and likely a direct driver
of, the dialect skew reported in \S\ref{subsec:dialect} (Egyptian and
Gulf Arabic dominant; Levantine and Maghrebi Arabic together under
7\%). Second, posting volume is heavily skewed toward recent years,
with negligible volume before 2022 and a sharp rise from 2024 onward,
reflecting both the underlying growth of these subreddits and the
collection window of this study; users and language patterns from
earlier years are correspondingly under-represented.

\section{Full Classification Results}
Tables \ref{tab:results_classical_full}--\ref{tab:results_classical_full2} (classical classifiers, split across two tables for length) and \ref{tab:results_plm_full} (PLMs) show the entire results from our two classification experiments.

\label{app:full_results}

\begin{table*}[t]
\centering
\footnotesize
\setlength{\tabcolsep}{4pt}
\caption{Full classical classifier results on CARMA (Reddit), conditions 1--6 of 12 (continued in Table~\ref{tab:results_classical_full2}). Each cell reports \textit{post level $|$ user level}: post-level instances treat each post as an independent, author-labeled example, while user-level instances aggregate a user's posts into a single TF-IDF document. User-level test sets are small for rare conditions; these F1 values are correspondingly higher-variance than the post-level results.}
\label{tab:results_classical_full}
\begin{tabular}{llcccc}
\toprule
\textbf{Condition} & \textbf{Model} & \textbf{F1} & \textbf{Acc.} & \textbf{P} & \textbf{R} \\
\midrule
\multirow{4}{*}{Depression}
 & NB & 0.807|0.488 & 0.679|0.606 & 0.681|0.697 & 0.992|0.375 \\
 & LR & 0.791|0.676 & 0.667|0.677 & 0.690|0.678 & 0.926|0.674 \\
 & SVM & 0.787|0.674 & 0.664|0.663 & 0.690|0.653 & 0.917|0.696 \\
 & XGB & \textbf{0.808}|\textbf{0.683} & 0.680|0.682 & 0.681|0.681 & 0.994|0.685 \\
\midrule
\multirow{4}{*}{Anxiety}
 & NB & 0.819|0.661 & 0.696|0.669 & 0.714|0.677 & 0.960|0.646 \\
 & LR & 0.798|0.744 & 0.676|0.761 & 0.721|0.804 & 0.894|0.692 \\
 & SVM & 0.782|0.797 & 0.659|0.808 & 0.721|0.845 & 0.853|0.754 \\
 & XGB & \textbf{0.828}|\textbf{0.850} & 0.709|0.862 & 0.717|0.927 & 0.979|0.785 \\
\midrule
\multirow{4}{*}{ADHD}
 & NB & 0.738|0.576 & 0.590|0.615 & 0.602|0.641 & 0.951|0.523 \\
 & LR & 0.733|0.667 & 0.592|0.654 & 0.608|0.643 & 0.921|0.692 \\
 & SVM & 0.724|0.662 & 0.590|0.646 & 0.612|0.634 & 0.885|0.692 \\
 & XGB & \textbf{0.751}|\textbf{0.721} & 0.603|0.685 & 0.607|0.646 & 0.984|0.815 \\
\midrule
\multirow{4}{*}{Sleep dis.}
 & NB & \textbf{0.810}|0.765 & 0.684|0.778 & 0.686|0.812 & 0.988|0.722 \\
 & LR & 0.781|0.754 & 0.655|0.764 & 0.688|0.788 & 0.903|0.722 \\
 & SVM & 0.752|0.771 & 0.629|0.778 & 0.690|0.794 & 0.827|0.750 \\
 & XGB & 0.802|\textbf{0.886} & 0.671|0.889 & 0.680|0.912 & 0.976|0.861 \\
\midrule
\multirow{4}{*}{OCD}
 & NB & 0.746|0.594 & 0.612|0.594 & 0.716|0.594 & 0.778|0.594 \\
 & LR & 0.716|0.739 & 0.597|0.734 & 0.740|0.727 & 0.692|0.750 \\
 & SVM & 0.683|0.742 & 0.570|0.750 & 0.742|0.767 & 0.633|0.719 \\
 & XGB & \textbf{0.797}|\textbf{0.821} & 0.672|0.844 & 0.728|0.958 & 0.882|0.719 \\
\midrule
\multirow{4}{*}{Panic}
 & NB & \textbf{0.825}|0.722 & 0.734|0.667 & 0.767|0.619 & 0.892|0.867 \\
 & LR & 0.824|0.812 & 0.727|0.800 & 0.754|0.765 & 0.909|0.867 \\
 & SVM & 0.804|\textbf{0.909} & 0.711|0.900 & 0.770|0.833 & 0.841|1.000 \\
 & XGB & 0.818|0.788 & 0.717|0.767 & 0.747|0.722 & 0.903|0.867 \\
\bottomrule
\end{tabular}
\end{table*}

\begin{table*}[t]
\centering
\footnotesize
\setlength{\tabcolsep}{4pt}
\caption{Full classical classifier results on CARMA (Reddit), conditions 7--12 of 12 (continued from Table~\ref{tab:results_classical_full}). Each cell reports \textit{post level $|$ user level}. User-level test sets are small for rare conditions (e.g., $n{=}10$ for eating disorder and schizophrenia, $n{=}16$ for bipolar disorder and suicidal ideation); these F1 values are correspondingly higher-variance than the post-level results.}
\label{tab:results_classical_full2}
\begin{tabular}{llcccc}
\toprule
\textbf{Condition} & \textbf{Model} & \textbf{F1} & \textbf{Acc.} & \textbf{P} & \textbf{R} \\
\midrule
\multirow{4}{*}{BPD}
 & NB & \textbf{0.863}|\textbf{0.720} & 0.778|0.650 & 0.808|0.600 & 0.927|0.900 \\
 & LR & 0.825|0.556 & 0.729|0.600 & 0.804|0.625 & 0.849|0.500 \\
 & SVM & 0.765|0.600 & 0.661|0.600 & 0.802|0.600 & 0.731|0.600 \\
 & XGB & 0.858|0.571 & 0.769|0.550 & 0.799|0.545 & 0.927|0.600 \\
\midrule
\multirow{4}{*}{PTSD}
 & NB & \textbf{0.665}|\textbf{0.700} & 0.550|0.700 & 0.514|0.700 & 0.944|0.700 \\
 & LR & 0.665|0.444 & 0.567|0.500 & 0.525|0.500 & 0.906|0.400 \\
 & SVM & 0.632|0.667 & 0.559|0.650 & 0.523|0.636 & 0.798|0.700 \\
 & XGB & 0.655|0.500 & 0.539|0.500 & 0.507|0.500 & 0.924|0.500 \\
\midrule
\multirow{4}{*}{Suicidal}
 & NB & 0.861|0.667 & 0.772|0.562 & 0.786|0.538 & 0.953|0.875 \\
 & LR & 0.855|0.706 & 0.764|0.688 & 0.786|0.667 & 0.936|0.750 \\
 & SVM & 0.813|\textbf{0.824} & 0.715|0.812 & 0.792|0.778 & 0.835|0.875 \\
 & XGB & \textbf{0.864}|0.750 & 0.776|0.750 & 0.784|0.750 & 0.963|0.750 \\
\midrule
\multirow{4}{*}{Bipolar}
 & NB & \textbf{0.883}|0.571 & 0.791|0.438 & 0.792|0.462 & 0.998|0.750 \\
 & LR & 0.874|0.526 & 0.778|0.438 & 0.793|0.455 & 0.975|0.625 \\
 & SVM & 0.803|0.556 & 0.682|0.500 & 0.790|0.500 & 0.816|0.625 \\
 & XGB & 0.871|\textbf{0.588} & 0.772|0.562 & 0.789|0.556 & 0.971|0.625 \\
\midrule
\multirow{4}{*}{Schizophrenia}
 & NB & \textbf{0.956}|0.769 & 0.915|0.700 & 0.917|0.625 & 0.998|1.000 \\
 & LR & 0.949|\textbf{0.909} & 0.903|0.900 & 0.922|0.833 & 0.977|1.000 \\
 & SVM & 0.875|0.833 & 0.782|0.800 & 0.919|0.714 & 0.835|1.000 \\
 & XGB & 0.943|0.600 & 0.892|0.600 & 0.915|0.600 & 0.972|0.600 \\
\midrule
\multirow{4}{*}{Eating dis.}
 & NB & 0.773|0.500 & 0.651|0.400 & 0.801|0.429 & 0.746|0.600 \\
 & LR & 0.761|0.667 & 0.642|0.600 & 0.811|0.571 & 0.716|0.800 \\
 & SVM & 0.745|\textbf{0.727} & 0.628|0.700 & 0.819|0.667 & 0.683|0.800 \\
 & XGB & \textbf{0.789}|0.500 & 0.667|0.400 & 0.795|0.429 & 0.783|0.600 \\
\bottomrule
\end{tabular}
\end{table*}

\begin{table*}[t]
\centering
\footnotesize
\setlength{\tabcolsep}{4pt}
\caption{Full fine-tuned PLM results on CARMA (Reddit). Each cell reports \textit{post level $|$ user level}: post-level instances classify each post individually (label inherited from its author), while user-level instances concatenate a user's posts into a single (truncated) input sequence. Best F1 per condition in \textbf{bold} for each level. User-level test sets are small for rare conditions (e.g., $n{=}10$ for eating disorder and schizophrenia, $n{=}16$ for bipolar disorder and suicidal ideation); these F1 values are correspondingly higher-variance than the post-level results.}
\label{tab:results_plm_full}
\begin{tabular}{llcccc}
\toprule
\textbf{Condition} & \textbf{Model} & \textbf{F1} & \textbf{Acc.} & \textbf{P} & \textbf{R} \\
\midrule
\multirow{2}{*}{Depression}
 & AraBERT & \textbf{0.796}|\textbf{0.707} & 0.680|0.674 & 0.702|0.642 & 0.918|0.788 \\
 & CAMeLBERT & 0.789|0.606 & 0.670|0.647 & 0.697|0.685 & 0.908|0.543 \\
\midrule
\multirow{2}{*}{Anxiety}
 & AraBERT & 0.766|0.696 & 0.643|0.623 & 0.722|0.583 & 0.817|0.862 \\
 & CAMeLBERT & \textbf{0.800}|\textbf{0.719} & 0.678|0.700 & 0.720|0.676 & 0.901|0.769 \\
\midrule
\multirow{2}{*}{ADHD}
 & AraBERT & 0.747|0.672 & 0.622|0.661 & 0.629|0.652 & 0.920|0.692 \\
 & CAMeLBERT & \textbf{0.749}|\textbf{0.710} & 0.629|0.692 & 0.635|0.671 & 0.914|0.754 \\
\midrule
\multirow{2}{*}{Sleep dis.}
 & AraBERT & \textbf{0.777}|0.719 & 0.655|0.653 & 0.695|0.604 & 0.881|0.889 \\
 & CAMeLBERT & 0.773|\textbf{0.727} & 0.652|0.708 & 0.697|0.683 & 0.867|0.778 \\
\midrule
\multirow{2}{*}{OCD}
 & AraBERT & 0.709|0.571 & 0.605|0.625 & 0.771|0.667 & 0.656|0.500 \\
 & CAMeLBERT & \textbf{0.711}|\textbf{0.721} & 0.594|0.625 & 0.744|0.574 & 0.680|0.969 \\
\midrule
\multirow{2}{*}{Panic}
 & AraBERT & \textbf{0.851}|0.750 & 0.777|0.733 & 0.803|0.706 & 0.906|0.800 \\
 & CAMeLBERT & 0.848|\textbf{0.774} & 0.773|0.767 & 0.802|0.750 & 0.898|0.800 \\
\midrule
\multirow{2}{*}{BPD}
 & AraBERT & \textbf{0.847}|0.696 & 0.766|0.650 & 0.836|0.615 & 0.858|0.800 \\
 & CAMeLBERT & 0.801|\textbf{0.714} & 0.715|0.600 & 0.845|0.556 & 0.762|1.000 \\
\midrule
\multirow{2}{*}{PTSD}
 & AraBERT & 0.677|0.588 & 0.675|0.650 & 0.639|0.714 & 0.719|0.500 \\
 & CAMeLBERT & \textbf{0.701}|\textbf{0.640} & 0.649|0.550 & 0.588|0.533 & 0.867|0.800 \\
\midrule
\multirow{2}{*}{Suicidal}
 & AraBERT & \textbf{0.866}|\textbf{0.778} & 0.788|0.750 & 0.816|0.700 & 0.922|0.875 \\
 & CAMeLBERT & 0.837|0.750 & 0.754|0.750 & 0.820|0.750 & 0.856|0.750 \\
\midrule
\multirow{2}{*}{Bipolar}
 & AraBERT & 0.841|\textbf{0.714} & 0.734|0.750 & 0.801|0.833 & 0.884|0.625 \\
 & CAMeLBERT & \textbf{0.858}|0.571 & 0.754|0.438 & 0.789|0.462 & 0.940|0.750 \\
\midrule
\multirow{2}{*}{Schizophrenia}
 & AraBERT & \textbf{0.875}|\textbf{0.800} & 0.781|0.800 & 0.917|0.800 & 0.837|0.800 \\
 & CAMeLBERT & 0.869|0.800 & 0.770|0.800 & 0.908|0.800 & 0.833|0.800 \\
\midrule
\multirow{2}{*}{Eating dis.}
 & AraBERT & \textbf{0.748}|\textbf{0.727} & 0.614|0.700 & 0.776|0.667 & 0.722|0.800 \\
 & CAMeLBERT & 0.720|0.667 & 0.585|0.500 & 0.775|0.500 & 0.672|1.000 \\
\bottomrule
\end{tabular}
\end{table*}

\section{Condition Keyword Lists}
\label{app:keywords}

Table~\ref{tab:keywords} lists the Arabic and English condition
keywords used by the pipeline for each of the 12 retained conditions.
Keywords marked with $\dagger$ are \textit{restricted}: they are
permitted only in combination with suffering-class phrases
(\foreignlanguage{arabic}{أعاني من}, \foreignlanguage{arabic}{بعاني من})
to reduce standalone false positives. Orthographic variants (hamza /
hamza-less; taa marbuta \foreignlanguage{arabic}{ة} /
\foreignlanguage{arabic}{ه}) are generated automatically and are not
listed separately.

\begin{table*}[t]
\centering
\small
\caption{Condition keyword inventory, 341 keyword forms including English and Arabic variants, across the 15 conditions the extraction pipeline was run on. Of these, 12 are retained in the released corpus (Table~\ref{tab:corpus_summary}); DID, Trichotillomania, and Paranoid PD yielded too few confirmed users and are excluded from the final corpus and downstream experiments, but are included here for completeness of the pattern description. $\dagger$ = restricted to suffering phrases only.}
\label{tab:keywords}
\begin{tabular}{p{2.8cm} p{2.2cm} p{10.2cm}}
\toprule
\textbf{Condition} & \textbf{Arabic Label} & \textbf{Representative Keywords (Arabic / English)} \\
\midrule
Schizophrenia &
\foreignlanguage{arabic}{الفصام} &
\foreignlanguage{arabic}{فصام، الفصام، شيزوفرينيا، اضطراب الفصام، انفصام الشخصية، الاضطراب الفصامي الوجداني، سكيزوأفكتيف} /
\textit{schizophrenia, schizoaffective disorder} \\
\addlinespace
Bipolar Disorder &
\foreignlanguage{arabic}{ثنائي القطب} &
\foreignlanguage{arabic}{ثنائي القطب، اضطراب ثنائي القطب، بايبولار، هوس خفيف، مانيا، سيكلوثيميا} /
\textit{bipolar disorder, bipolar I/II, hypomania, mania, cyclothymia} \\
\addlinespace
Depression &
\foreignlanguage{arabic}{الاكتئاب} &
\foreignlanguage{arabic}{اكتئاب، الاكتئاب، اكتئاب حاد $/$ مزمن، اكتئاب ما بعد الولادة، الاكتئاب الموسمي، ديبريشن، سوء المزاج} /
\textit{depression, MDD, dysthymia, postpartum depression, seasonal affective disorder} \\
\addlinespace
Anxiety &
\foreignlanguage{arabic}{اضطراب القلق} &
\foreignlanguage{arabic}{اضطراب القلق، قلق عام $/$ مرضي $/$ مزمن، رهاب اجتماعي، قلق اجتماعي، رهاب الخلاء، قلق الانفصال} /
\textit{generalized anxiety, social anxiety, GAD, agoraphobia, specific phobia} \\
\addlinespace
Panic Disorder &
\foreignlanguage{arabic}{اضطراب الهلع} &
\foreignlanguage{arabic}{نوبة هلع، نوبات هلع، الهلع، اضطراب الهلع، نوبة ذعر} /
\textit{panic attack, panic disorder} \\
\addlinespace
OCD &
\foreignlanguage{arabic}{الوسواس القهري} &
\foreignlanguage{arabic}{وسواس قهري، الوسواس القهري، الوسواس الديني، وسواس الطهارة، ديسمورفيا، اكتناز مرضي} /
\textit{OCD, body dysmorphic disorder, BDD, hoarding disorder, scrupulosity} \\
\addlinespace
PTSD &
\foreignlanguage{arabic}{اضطراب ما بعد الصدمة} &
\foreignlanguage{arabic}{اضطراب ما بعد الصدمة، صدمة نفسية، صدمة معقدة، كرب ما بعد الصدمة} /
\textit{PTSD, C-PTSD, complex PTSD, acute stress disorder} \\
\addlinespace
Eating Disorder &
\foreignlanguage{arabic}{اضطراب الأكل} &
\foreignlanguage{arabic}{فقدان الشهية العصبي، أنوركسيا، الشره المرضي، بوليميا، نهم الطعام، $ARFID$} /
\textit{anorexia, bulimia, binge eating disorder, ARFID, orthorexia} \\
\addlinespace
ADHD &
\foreignlanguage{arabic}{فرط الحركة وتشتت الانتباه} &
\foreignlanguage{arabic}{$ADHD$، فرط حركة، تشتت انتباه، نقص الانتباه، اضطراب الانتباه} /
\textit{ADHD, attention deficit hyperactivity disorder; ADD$\dagger$} \\
\addlinespace
DID &
\foreignlanguage{arabic}{اضطراب الهوية التفارقي} &
\foreignlanguage{arabic}{اضطراب الهوية التفارقي، تعدد الشخصيات} /
\textit{dissociative identity disorder, DID, multiple personality disorder} \\
\addlinespace
Sleep Disorder &
\foreignlanguage{arabic}{اضطراب النوم} &
\foreignlanguage{arabic}{الأرق، أرق، انقطاع النفس النومي، فرط النوم، الخدار} /
\textit{insomnia, sleep apnea, hypersomnia, narcolepsy} \\
\addlinespace
Trichotillomania &
\foreignlanguage{arabic}{اضطراب شد الشعر} &
\foreignlanguage{arabic}{شد الشعر، نتف الشعر، نتف الجلد، قضم الأظافر} /
\textit{trichotillomania, excoriation disorder, onychophagia} \\
\addlinespace
BPD &
\foreignlanguage{arabic}{اضطراب الشخصية الحدية} &
\foreignlanguage{arabic}{اضطراب الشخصية الحدية، الحدية، بوردرلاين} /
\textit{borderline personality disorder, BPD} \\
\addlinespace
Paranoid PD &
\foreignlanguage{arabic}{اضطراب الشخصية الزوراني} &
\foreignlanguage{arabic}{اضطراب الشخصية الزوراني، الشخصية البارانوية، البارانويا} /
\textit{paranoid personality disorder} \\
\addlinespace
Suicidal Ideation &
\foreignlanguage{arabic}{الأفكار الانتحارية} &
\foreignlanguage{arabic}{أفكار انتحارية، الفكر الانتحاري} /
\textit{suicidal thoughts, suicidal ideation;
\foreignlanguage{arabic}{انتحار}
$\dagger$، \foreignlanguage{arabic}{محاولة انتحار}$\dagger$، \foreignlanguage{arabic}{إيذاء النفس}$\dagger$ / suicide attempt$\dagger$, self-harm$\dagger$} \\
\bottomrule
\end{tabular}
\end{table*}

\section{Self-Disclosure Phrase Inventory}
\label{app:phrases}

Tables~\ref{tab:phrases_universal} and~\ref{tab:phrases_dialect} list
the full self-disclosure phrase inventory (117 universal + 30 dialect-specific not in the universal set = 147 total unique phrases). Universal phrases
(Table~\ref{tab:phrases_universal}) apply across all dialects; they are
organized by semantic type and connector class. Dialect-specific phrases
(Table~\ref{tab:phrases_dialect}) supplement the universal set with
register-appropriate variants for Gulf, Levantine, Egyptian, Maghrebi,
and Modern Standard Arabic (MSA). Orthographic variants (hamza /
hamza-less; taa marbuta \foreignlanguage{arabic}{ة} /
\foreignlanguage{arabic}{ه}) are generated automatically at runtime and
not listed separately.

\begin{table}[t]
\centering
\small
\caption{Universal self-disclosure phrases by type and connector class.
\textbf{B} = \foreignlanguage{arabic}{بـ} connector;
\textbf{M} = \foreignlanguage{arabic}{من} connector;
\textbf{D} = direct (no connector).}
\label{tab:phrases_universal}
\resizebox{\columnwidth}{!}{%
\begin{tabular}{p{2.6cm} c p{5.6cm}}
\toprule
\textbf{Type} & \textbf{Conn.} & \textbf{Phrases} \\
\midrule
Affliction & B &
\foreignlanguage{arabic}{أنا مصاب $/$ انا مصاب، أنا مصابة، أنا مريض، أنا مريضة، كشخص مصاب} \\
\addlinespace
Suffering & M &
\foreignlanguage{arabic}{أعاني من $/$ اعاني من، أنا أعاني من، أنا أتعالج من، أنا أعيش مع، أصبت} \\
\addlinespace
Clinical Diagnosis & B &
\foreignlanguage{arabic}{تم تشخيصي، شخصوني، شخصني، شخصتني، شخصت، شخصني الطبيب، شخصني الدكتور، شخصتني الطبيبة، شخصتني الدكتورة، الأخصائي شخصني، الأخصائية شخصتني} \\
\addlinespace
Possessive & D &
\foreignlanguage{arabic}{عندي، أنا عندي} \\
\addlinespace
Discovery / & B &
\foreignlanguage{arabic}{طلعت مصاب، لقيت نفسي مصاب، اكتشفت إني مصاب، عرفت إن أنا مصاب، اتأكدت إني مصاب، كنت مصاب} \\
Confirmation & D &
\foreignlanguage{arabic}{تشخيصي هو، تشخيصي كان، طلع تشخيصي، طلعت عندي، لقيت نفسي عندي، اكتشفت إن عندي، عرفت إني عندي، اتأكدت إن عندي، بعد التشخيص طلع عندي، التحليل أثبت إن عندي} \\
\addlinespace
Doctor Result & B &
\foreignlanguage{arabic}{أكد الطبيب إصابتي، الدكتور أكد إصابتي، تم التأكيد على إصابتي، تم إثبات إصابتي، الدكتور أكد التشخيص، الطبيب أثبت التشخيص} \\
& D &
\foreignlanguage{arabic}{تشخيص الطبيب هو، تشخيص الدكتور هو، تشخيص حالتي هو، تشخيص المرض هو، تم تشخيص حالتي على أنها، التشخيص الذي حصلت عليه هو، طلع إن التشخيص هو، اتضح إن التشخيص كان} \\
\bottomrule
\end{tabular}
}
\end{table}

\begin{table}[t]
\centering
\small
\caption{Dialect-specific self-disclosure phrases.
\textbf{B} = \foreignlanguage{arabic}{بـ} connector;
\textbf{M} = \foreignlanguage{arabic}{من} connector;
\textbf{D} = direct.}
\label{tab:phrases_dialect}
\resizebox{\columnwidth}{!}{%
\begin{tabular}{p{1.5cm} c p{6.2cm}}
\toprule
\textbf{Dialect} & \textbf{Conn.} & \textbf{Phrases} \\
\midrule
Gulf &
B & \foreignlanguage{arabic}{انا مصاب، مبتلي، مبتلية، ابتليت} \\
& D & \foreignlanguage{arabic}{عندي، فيني، انا متعايش، جاني، جاتني} \\
& M & \foreignlanguage{arabic}{اعاني من} \\
\addlinespace
Levantine &
B & \foreignlanguage{arabic}{انا مصاب} \\
& D & \foreignlanguage{arabic}{عندي، معي، طلعلي، طلعتلي، صار معي، صارت معي، اجاني، اجاتني} \\
& M & \foreignlanguage{arabic}{بعاني من، عم بعاني} \\
\addlinespace
Egyptian &
B & \foreignlanguage{arabic}{انا مصاب، طلعت مصاب، لقيت نفسي مصاب، اكتشفت إني مصاب، اتأكدت إني مصاب} \\
& D & \foreignlanguage{arabic}{عندي، معايا، جالي، جاتلي، طلعلي، طلعت عندي، لقيت نفسي عندي، اكتشفت إن عندي، اتأكدت إن عندي، عرفت إني عندي، أنا بواجه مشكلة اسمها} \\
& M & \foreignlanguage{arabic}{بعاني من، اتشخصت} \\
\addlinespace
Maghrebi &
B & \foreignlanguage{arabic}{انا مصاب} \\
& D & \foreignlanguage{arabic}{عندي، كاين عندي، جاني، جاتني، طلعلي، تشخصت} \\
\addlinespace
MSA &
B & \foreignlanguage{arabic}{أنا مصاب، أنا مصابة، شُخِّصت، شخصت} \\
& D & \foreignlanguage{arabic}{لدي، تشخيصي هو، طلع تشخيصي، أتعايش مع} \\
& M & \foreignlanguage{arabic}{أعاني من، تم تشخيصي} \\
\bottomrule
\end{tabular}
}
\end{table}

\section{Connector Logic and Orthographic Variants}
\label{app:connectors}

Table~\ref{tab:connectors} summarizes the three connector types used by
the pipeline, showing how each phrase class links to the condition
keyword and how the definite article \foreignlanguage{arabic}{ال}
triggers the \foreignlanguage{arabic}{بال} fusion rule.

\begin{table}[h]
\centering
\small
\caption{Connector rules: phrase type, connector morpheme, keyword form generated, and example.}
\label{tab:connectors}
\resizebox{\columnwidth}{!}{%
\begin{tabular}{p{2.0cm} c p{3.0cm} p{3.5cm}}
\toprule
\textbf{Phrase Type} & \textbf{Conn.} & \textbf{Keyword Form} & \textbf{Example} \\
\midrule
Affliction, Diagnosis, Confirmation &
\foreignlanguage{arabic}{بـ} &
Base kw: \foreignlanguage{arabic}{بـ}+kw \newline
\foreignlanguage{arabic}{ال}-kw: \foreignlanguage{arabic}{بال}+stem \newline
English: \foreignlanguage{arabic}{بـ} kw &
\foreignlanguage{arabic}{مصاب باكتئاب} \newline
\foreignlanguage{arabic}{مصاب بالاكتئاب} \newline
\foreignlanguage{arabic}{مصاب بـ $PTSD$} \\
\addlinespace
Suffering &
\foreignlanguage{arabic}{من} &
Phrase ends in \foreignlanguage{arabic}{من}: bare kw \newline
Phrase lacks \foreignlanguage{arabic}{من}: \foreignlanguage{arabic}{من}+kw &
\foreignlanguage{arabic}{أعاني من الاكتئاب} \newline
\foreignlanguage{arabic}{عم بعاني من اكتئاب} \\
\addlinespace
Possessive, Discovery &
Direct &
Bare or definite kw &
\foreignlanguage{arabic}{عندي اكتئاب} \newline
\foreignlanguage{arabic}{تشخيصي هو الفصام} \\
\bottomrule
\end{tabular}
}
\end{table}

\noindent The pipeline generates all orthographic variants automatically.
Table~\ref{tab:ortho} lists the variant pairs handled.

\begin{table}[h]
\centering
\small
\caption{Orthographic variant pairs generated automatically.}
\label{tab:ortho}
\begin{tabular}{lll}
\toprule
\textbf{Feature} & \textbf{Variant A} & \textbf{Variant B} \\
\midrule
Hamza on alef & \foreignlanguage{arabic}{أنا} & \foreignlanguage{arabic}{انا} \\
Hamza in conjunction & \foreignlanguage{arabic}{إن} & \foreignlanguage{arabic}{ان} \\
Taa marbuta & \foreignlanguage{arabic}{مصابة} & \foreignlanguage{arabic}{مصابه} \\
Taa marbuta & \foreignlanguage{arabic}{مريضة} & \foreignlanguage{arabic}{مريضه} \\
Hamza on \foreignlanguage{arabic}{ا} & \foreignlanguage{arabic}{أعاني} & \foreignlanguage{arabic}{اعاني} \\
\bottomrule
\end{tabular}
\end{table}


\section{Annotation Guidelines}
\label{app:guidelines}

Both human annotators received the following written guidelines
(translated from Arabic and lightly edited for clarity, but otherwise
reproduced in full) prior to labeling. Annotators judged each post
together with its pipeline-assigned \textit{diagnosis} column and
assigned exactly one of two labels: \texttt{TP} or \texttt{FP}.

\paragraph{True Positive (TP).} Based on the post content, does the
annotator believe this person genuinely has the mental health condition
named in the \textit{diagnosis} column ,  either because the pipeline
captured a genuine disclosure, or because the post describes matching
symptoms or medications? An explicit phrase such as ``I was officially
diagnosed'' is \textit{not} required.

\begin{quote}
\foreignlanguage{arabic}{انا عندي اكتئاب} ,  \textit{``I have depression.''}\\
\foreignlanguage{arabic}{اتشخصت بـ اكتئاب} ,  \textit{``I was diagnosed with depression.''}
\end{quote}

\paragraph{False Positive (FP).} A post whose phrasing matches the
pipeline's keyword pattern but should not be labeled a genuine
disclosure because of a shortcoming in the pattern match, such as
sarcasm or reference to a third party.

\begin{quote}
\foreignlanguage{arabic}{اخويا عنده اكتئاب} ,  \textit{``My brother has depression.''} (third-person reference)\\
\foreignlanguage{arabic}{الاكتئاب هيموتني ههههههه} ,  \textit{``Depression is going to kill me, hahaha.''} (sarcasm)
\end{quote}

\section{Agreement Computation Details}
\label{app:agreement_details}

\begin{table}[h]
\centering
\small
\begin{tabular}{lcc}
\toprule
 & \textbf{Ann-2: TP} & \textbf{Ann-2: FP} \\
\midrule
\textbf{Ann-1: TP} & 172 & 6 \\
\textbf{Ann-1: FP} & 18 & 39 \\
\bottomrule
\end{tabular}
\caption{Confusion matrix between the two human annotators (Ann-1,
Ann-2) on the 235 posts both annotators labeled.}
\label{tab:iaa_confusion}
\end{table}


\subsection*{Kappa interpretation scale}

Cohen's $\kappa$ values are interpreted following
\citet{landis1977measurement}: $\kappa \leq 0.20$ slight,
$0.21$--$0.40$ fair, $0.41$--$0.60$ moderate, $0.61$--$0.80$
substantial, $0.81$--$1.00$ almost perfect agreement.

\subsection*{LLM Judges as Additional Annotators}
\label{app:llm_fp_estimate}

As a robustness check, we additionally had two LLM judges, Qwen2-7B-Instruct \citep{yang-etal-2024-qwen2} and Jais-2-8B-Chat \citep{anwar-etal-2025-jais2} (Qwen and Jais below), independently label the same 240-post sample used for human
annotation (\S\ref{subsec:human_annotation}) under the identical
\texttt{TP}/\texttt{FP} scheme. Table~\ref{tab:iaa_full} reports
pairwise agreement across all four raters.

\begin{table}[h]
\centering
\small
\begin{tabular}{lccc}
\toprule
\textbf{Pair} & \textbf{n} & \textbf{\% agree} & \textbf{$\kappa$} \\
\midrule
Ann-1 vs.\ Ann-2 & 235 & 89.8 & 0.701 \\
Ann-1 vs.\ Qwen & 229 & 76.4 & 0.167 \\
Ann-1 vs.\ Jais & 227 & 74.9 & 0.171 \\
Ann-2 vs.\ Qwen & 224 & 79.0 & 0.138 \\
Ann-2 vs.\ Jais & 224 & 77.7 & 0.176 \\
Qwen vs.\ Jais & 218 & 83.0 & 0.231 \\
\bottomrule
\end{tabular}
\caption{Pairwise agreement among all four raters (two human
annotators, two LLM judges) on the 240-post sample. Per Landis and
Koch's (1977) scale, human-human agreement is \emph{substantial}
($\kappa=0.701$) while every pair involving an LLM judge is only
\emph{slight-to-fair} ($\kappa = 0.138$--$0.231$).}
\label{tab:iaa_full}
\end{table}

Each rater's individual \texttt{FP} rate, averaged over the pairwise
subsets in which that rater participates, is: Ann-1 23.2\%, Ann-2 18.1\%,
Qwen 11.2\%, Jais 13.8\%. Two pooled estimates follow from these rates:
averaging all four raters gives 16.6\%, while averaging the two LLM
judges alone gives 12.5\% -- both substantially lower than the
21.7\% human-only estimate reported in
\S\ref{subsec:human_annotation}.

We do not fold the LLM judges into the headline estimate, because the
gap above reflects a systematic bias rather than additional independent
noise around the same true rate. Conditioning on the posts each human
annotator labeled \texttt{FP}, the LLM judges recover only a minority
of them: of Ann-1's 51 (vs.\ Qwen) and 52 (vs.\ Jais) \texttt{FP} calls,
Qwen agrees on 21.6\% and Jais on 25.0\%; of Ann-2's 38 (vs.\ Qwen) and
41 (vs.\ Jais) \texttt{FP} calls, Qwen agrees on 21.1\% and Jais on
26.8\%. In other words, LLM judges default to \texttt{TP} on many posts
that both human annotators independently flag as sarcasm, third-person
reference, figurative use, or negation -- the same non-literal
constructions the pattern-based pipeline itself struggles with. Their
lower apparent FP rate is therefore evidence of lower sensitivity to
these cases, not evidence that the corpus is cleaner than the
human-only estimate suggests.

This pattern is consistent with three independent strands of prior
work. First, sarcasm and other forms of dialectal figurative language
are a documented weak point for Arabic NLP specifically, motivating
dedicated resources such as the ArSarcasm dataset
\citep{abu-farha-magdy-2020-arabic}. Second, negation handling is a
well-known blind spot for large language models more broadly, which
tend to default to the affirmative reading of a statement even when a
negation cue is present \citep{garcia-ferrero-etal-2023-not}. Third,
using LLMs as binary judges has independently been shown to induce a
general \emph{leniency} (or ``agreeableness'') bias -- high agreement
on positive cases alongside a much lower rate of correctly identifying
negative ones \citep{jain2025beyond}, mirroring the asymmetry we
observe here (Qwen and Jais both agree with humans far more often on
\texttt{TP} than on \texttt{FP}; Table~\ref{tab:iaa_full}). Together
these strands suggest our LLM judges' lower FP rate reflects known,
generic limitations of LLM-based annotation rather than a
sample-specific artifact.

\subsection*{Illustrative Disagreement Examples}

Table~\ref{tab:disagreement_examples} illustrates the four most common
disagreement patterns among the four raters (cf.\ the disagreement
breakdown underlying Table~\ref{tab:iaa_full}) with paraphrased
examples. To protect user privacy, consistent with our policy of
removing self-disclosure posts from the released corpus
(\S\ref{par:sensitive_data}), these examples are \textit{paraphrased
reconstructions} illustrating the phenomenon behind each disagreement
pattern rather than verbatim quotations from CARMA.

\begin{table*}[t]
\centering
\footnotesize
\renewcommand{\arraystretch}{1.15}
\begin{tabular}{@{}p{0.97\textwidth}@{}}
\toprule
\textbf{Paraphrased Disagreement Examples} \\
\midrule
\textit{``Depression is my old friend, we've known each other for ages
haha''} \newline
\textbf{Ann-1:} FP \quad \textbf{Ann-2:} FP \quad \textbf{Qwen:} TP
\quad \textbf{Jais:} TP \quad \textbf{Likely source:} sarcasm read
literally by LLMs \\
\addlinespace[4pt] \hline \addlinespace[4pt]
\textit{``I slept early last night, thank God insomnia didn't
come''} \newline
\textbf{Ann-1:} FP \quad \textbf{Ann-2:} TP \quad \textbf{Qwen:} TP
\quad \textbf{Jais:} FP \quad \textbf{Likely source:} negation missed
by Ann-2/Qwen \\
\addlinespace[4pt] \hline \addlinespace[4pt]
\textit{``My sister has been suffering from anxiety for a
while''} \newline
\textbf{Ann-1:} FP \quad \textbf{Ann-2:} FP \quad \textbf{Qwen:} TP
\quad \textbf{Jais:} TP \quad \textbf{Likely source:} third-person
subject conflated with poster \\
\addlinespace[4pt] \hline \addlinespace[4pt]
\textit{``I have some obsessiveness but I'm not sure if it's really
OCD''} \newline
\textbf{Ann-1:} FP \quad \textbf{Ann-2:} TP \quad \textbf{Qwen:} TP
\quad \textbf{Jais:} TP \quad \textbf{Likely source:} hedged/weak
self-diagnosis judged differently \\
\bottomrule
\end{tabular}
\caption{Paraphrased examples of the four most frequent multi-rater
disagreement patterns (analogous to the disagreement-pattern breakdown
produced by our agreement-analysis script). Each example is a
reconstructed illustration of the underlying phenomenon, not a verbatim
CARMA post.}
\label{tab:disagreement_examples}
\end{table*}

\end{document}
